\section*{TÓM TẮT ĐỒ ÁN}
\addcontentsline{toc}{section}{\numberline {}TÓM TẮT ĐỒ ÁN}
6G hướng tới phủ sóng toàn cầu bằng mạng phi mặt đất (NTN), trong đó vệ tinh quỹ đạo thấp (LEO) là thành phần then chốt. Do bay nhanh và vùng phủ hẹp, một thiết bị người dùng (UE) trong chòm vệ tinh LEO mật độ cao (kiểu Starlink) phải chuyển giao (handover) liên tục, trong khi việc đo tín hiệu kiểu mạng mặt đất không còn đáng tin cậy. Bài báo \textit{``Intelligent Handover Scheme for Improved 6G NTN LEO Satellite Network Performance''} (gọi tắt \textbf{ILCHO}) đề xuất kết hợp cơ chế Chuyển giao có điều kiện (Conditional Handover - CHO) của 3GPP với học tăng cường đa tác tử (QMIX) để chọn vệ tinh đích cho từng UE, dựa trên nền tảng giải tích chứng minh vị trí vệ tinh là tất định và bài toán gán tối ưu là NP-hard.

Trong đồ án này, em thực hiện hai phần việc chính. \textbf{Thứ nhất}, em khảo sát (survey) và phân tích chi tiết bài báo ILCHO: bối cảnh, mô hình hệ thống, thiết kế thuật toán, kết quả công bố, đồng thời chỉ ra các điểm mạnh, hạn chế và khoảng trống nghiên cứu khi đặt bài báo vào bức tranh chung của các phương pháp AI cho handover. \textbf{Thứ hai}, vì bài báo không công khai mã nguồn và không công bố đầy đủ tham số, em tự xây dựng lại toàn bộ mô hình (hình học chòm vệ tinh, mô hình kênh 3GPP, môi trường mô phỏng CHO đa tác tử, bộ điều khiển QMIX và bốn phương pháp so sánh) để tái hiện thực nghiệm, huấn luyện qua ba lượt tăng dần quy mô tới đúng quy mô bài báo công bố (100 UE, 600 giây/lượt phục vụ), và kiểm chứng độ ổn định qua nhiều lần huấn luyện độc lập (seed).

Kết quả tái hiện cho thấy tuyên bố cốt lõi của bài báo — cơ chế CHO điều khiển bằng học tăng cường đa tác tử giữ số chuyển giao thất bại thấp và ổn định trong chòm vệ tinh mật độ cao, trong khi các phương pháp truyền thống (chọn vệ tinh gần nhất, chọn vệ tinh có thời gian nhìn thấy lâu nhất) sụp đổ khi tải tăng — được tái hiện rõ ràng và vững chắc, đặc biệt khi huấn luyện đúng quy mô bài báo. Đồng thời, em cũng phát hiện và báo cáo trung thực hai điểm không tái hiện được (thông lượng cao nhất; sự bất ổn của biến thể hàm thưởng tuyến tính), cùng các hạn chế còn tồn tại do một số tham số quan trọng không được bài báo công bố.

\vspace{6pt}
\noindent\textbf{Từ khóa:} 6G, mạng phi mặt đất (NTN), vệ tinh LEO, chuyển giao có điều kiện (CHO), học tăng cường đa tác tử (MARL), QMIX, ILCHO.
